\documentclass[../main.tex]{subfiles}

\begin{document}
\chapter{Conclusioni}
\label{chap:conclusion}

Alla base di questo lavoro vi è un quesito fondamentale: \textit{quanto contano i dati, rispetto all'architettura, nelle performance di un modello deep learning?} I risultati ottenuti non danno una risposta unica, ma una più interessante: \textbf{dipende da quale errore, e da quale modello}.

Il rumore sulle etichette emerge come il fattore di criticità più rilevante, rispetto al quale né le proprietà strutturali del dataset, né la capacità di generalizzazione temporale, né eventuali meccanismi di fallback risultano sufficienti ad attenuarne gli effetti. Quando il segnale di supervisione è corrotto, entrambi i modelli subiscono un netto degrado prestazionale, con un impatto ancora più marcato nel caso di LSTM rispetto a MLP. Questo risultato è particolarmente significativo, poiché suggerisce che la robustezza architetturale presenta un limite ben definito, collocato nel processo di etichettatura. In un contesto industriale reale, nel quale le etichette derivano spesso da log di manutenzione compilati manualmente o da soglie di allarme definite in modo approssimativo, tale evidenza assume implicazioni operative dirette: investire nella qualità delle etichette risulta più efficace che intervenire esclusivamente sul raffinamento dell'architettura~\cite{jakubik2024datacentricartificialintelligence}.

Il comportamento sotto corruzione delle feature racconta invece una storia diversa per i due modelli, e questa è la lezione più originale del lavoro. I due modelli, a parità di dato in ingresso, sviluppano strategie di compensazione radicalmente diverse. MLP sfrutta la ridondanza statica tra feature correlate: quando \texttt{TP3} viene corrotta, si appoggia a \texttt{Reservoirs} quasi automaticamente, come se la correlazione perfetta tra le due colonne fosse una rete di sicurezza implicita. LSTM invece sfrutta la ridondanza temporale: compensa i valori corrotti di \texttt{Oil\_temperature} attingendo al contesto delle finestre adiacenti, ma su \texttt{TP3} questo meccanismo è meno efficace, perché la struttura sequenziale non sostituisce altrettanto bene la correlazione tra colonne. Architetture diverse imparano a fare affidamento su dimensioni diverse dell'informazione, e questo le rende vulnerabili a tipologie di errore diverse.

Una terza implicazione riguarda la metrica F1-score weighted, la quale, su un dataset con sbilanciamento 98\%--1.3\%, è quasi cieca alla corruzione delle feature, infatti il modello può perdere gran parte della propria capacità discriminante sulla classe di guasto senza che la metrica aggregata lo rilevi. Questo non rappresenta un limite specifico del solo dataset analizzato, ma una criticità metodologica più generale~\cite{buda2018systematic, johnson2019survey}: in presenza di dati sbilanciati, affidarsi a un'unica metrica di valutazione può condurre a conclusioni sistematicamente troppo ottimistiche sulla robustezza del modello.

\paragraph{Sguardo al futuro.}
I risultati di questo studio aprono almeno due direzioni che vale la pena esplorare. \textbf{La prima} è replicare l'analisi su dataset di domini diversi, per capire se le vulnerabilità osservate siano proprietà del dataset MetroPT-3 o pattern più generali. \textbf{La seconda} è testare pattern di corruzione strutturata e correlata nel tempo (più vicini ai guasti reali dei sensori industriali) in sostituzione del meccanismo MCAR adottato in questo lavoro.

Sul fronte dello sviluppo di PuckTrick, i risultati del metodo \textit{noise} su \texttt{DV\_pressure} suggeriscono un limite dell'attuale strategia di perturbazione: generare valori uniformi nel range della feature produce un rumore semanticamente insensato quando la distribuzione è dominata da pochi valori discreti. Un miglioramento naturale sarebbe introdurre una logica adattiva che, qualora due o più valori coprano più del 90\% del dataset, ignori i valori residui e campioni randomicamente tra quei valori dominanti, preservando la struttura informativa della feature anziché diluirla.

Un secondo limite riguarda la natura intrinsecamente tabulare dei metodi di corruzione attualmente implementati: PuckTrick tratta ogni riga in modo indipendente, senza alcuna consapevolezza dell'ordine temporale. Questo approccio, come discusso nel Capitolo~\ref{chap:StateOfArt}, produce esclusivamente \textit{point anomalies}, mentre nei sistemi reali le anomalie si manifestano spesso come \textit{contextual anomalies}: valori normali in assoluto ma anomali nel contesto della sequenza in cui compaiono. Estendere PuckTrick con metodi di corruzione time-aware, capaci di introdurre rumore che tenga conto della struttura temporale del dato, \textit{ad esempio burst di valori mancanti contigui, drift graduali di una feature, o inversioni localizzate in una finestra temporale}, consentirebbe di simulare scenari di degrado più realistici e di valutare la robustezza dei modelli sequenziali in condizioni più vicine a quelle operative.

Un terzo limite riguarda l'implementazione del rumore su colonne di tipo \textit{date} all'interno del metodo \textit{noise}. Come descritto nel capitolo~\ref{chap:DesignAndImplementation}, paragrafo~\ref{paragraph:eccezioneColonnaDate}, tale implementazione presenta criticità di ottimizzazione che possono incidere negativamente sulle prestazioni complessive del metodo.

Infine, un limite tecnico emerso durante l'integrazione riguarda la gestione dell'engine di esecuzione: attualmente, se l'utente seleziona il backend Spark ma fornisce un dataset non-Spark, la libreria non esegue la conversione automaticamente e non viene materializzato l'identificativo \textit{\_pucktrick\_row\_id} (Capitolo~\ref{subsection:problemaIDOriginalDF}), ma invece ritorna un errore parlante. Rendere questa conversione trasparente, rilevando il tipo del dataset in input e convertendolo automaticamente in un DataFrame Spark quando necessario, migliorerebbe significativamente l'esperienza d'uso e abbasserebbe la barriera di adozione della libreria in contesti distribuiti.
\end{document}